\section{R3': finite-horizon cone confinement}
\label{sec:R3prime}

The third result describes the geometry of the trajectory itself, not
only its endpoint: with high probability the angle to the correct-token
row enters and stays in a spherical cap, and the terminal angular margin
concentrates. This is a \emph{finite-horizon, high-probability}
confinement statement proved by the same martingale machinery as
\Cref{thm:R1,thm:R2}. It is \emph{not} a differential equation: there is
no drift ODE, no closability of a summary statistic, and no per-step
Foster--Lyapunov rate. The deterministic ODE limit appears only as a
clearly flagged informal remark (\Cref{rem:R3prime_ode}) under two extra
stylizations that are \emph{not} implied by \Cref{ass:descent}.

\begin{theorem}[Finite-horizon cone confinement]
\label{thm:R3prime}
Suppose \Cref{ass:descent,ass:bounded_architecture} hold with the mild
incoherence \Cref{eq:R1_incoherence} and $\snet\ge\sqrt2\,\critrate$ (the
super-critical regime of \Cref{thm:R1}). Fix a horizon $T_{\max}\ge1$, a
tolerance $\varepsilon\in(0,1)$, and $\delta_{\mathrm{fail}}\in(0,1)$. Let
$m^\star\coloneqq\snet\driftc\rho_0/M$ be the target angular level (the
signal floor of \Cref{lem:signal_floor} in angular-margin units, an $O(1)$
cosine), and let $C\coloneqq\{x\in\Sphere:\snorder(x)\ge m^\star-\varepsilon\}$
be the spherical cap of states whose angular margin to $W_U^{\corr}$ is at
least $m^\star-\varepsilon$. Then, on an event of probability at least
$1-\delta_{\mathrm{fail}}-(|\Vocab|-1)^{-1}$:
\begin{enumerate}[label=(\roman*),leftmargin=2em,itemsep=0.2em]
\item \emph{Entry and confinement.} There is a burn-in time
       $T_0=O(1)$ in rescaled time (equivalently
       $T_0=O(\sigma_{T_{\max}}^{-2}\,\varepsilon^{2})$ raw steps) such
       that $x_t\in C$ for every $t$ with $T_0\le t\le T_{\max}$.
\item \emph{Terminal concentration.} The terminal angular margin
       concentrates at the signal level,
       \begin{align}\label{eq:R3prime_concentration}
          \Bigl|\, \snorder_{T_{\max}} - m^\star \,\Bigr|
          \;\le\;
          O\!\left(\frac{1}{\sqrt{T_{\max}\,d}}\right),
       \end{align}
       with the same $1/\sqrt{T_{\max}d}$ fluctuation scale as
       \Cref{thm:R1,thm:R2}.
\end{enumerate}
\end{theorem}

\begin{proof}
The argument applies the signal floor and the martingale tail
\emph{uniformly over the horizon} via the maximal inequality already inside
\Cref{lem:azuma}, not via a drift recurrence. No ODE, no closability, no
per-step contraction is used. Abbreviate $T=T_{\max}$.

\smallskip\noindent\textbf{Union-bound paragraph.} The
$1-\delta_{\mathrm{fail}}-(|\Vocab|-1)^{-1}$ budget is the same good event
as \Cref{thm:R1}: the signal-floor event of \Cref{lem:signal_floor}
(budget $\delta_{\mathrm{fail}}$) intersected with the incorrect-max event
of \Cref{lem:incorrect_max} (budget $(|\Vocab|-1)^{-1}$), by a union bound
over these two events. The maximal-inequality form of \Cref{lem:azuma} (the
$\max_t$ inside \Cref{eq:azuma_tail}) supplies the \emph{simultaneous}
control over all $t\le T$ at no extra budget. Anti-pole non-degeneracy
$+$ $x_0$-washout absorbed as in \Cref{thm:R1} (no initial condition).

\smallskip\noindent\textbf{(i) Entry and confinement.} For each $t$, the
angular margin is $\snorder_t=\inner{\rowhat}{\tilde x_t}/\norm{\tilde x_t}_2$
(\Cref{def:angular_margin,lem:running_average}). By the signal-floor drift
bound (\Cref{lem:signal_floor} Step~2) the predictable part of
$\inner{\rowhat}{\tilde x_t}$ is $\ge\snet\driftc\rho_0$ for every $t$ once
the $x_0$-contribution $w_{t,0}=O(1/t)$ has decayed below $\varepsilon M/2$,
i.e.\ for $t\ge T_0$ with $T_0=O(1/\varepsilon)$ in the rescaled time
$\tau=1/t$ (\Cref{lem:running_average}); combined with $\norm{\tilde x_t}_2\le M$
this gives $\E[\snorder_t\mid\Fcal_0]\ge\snet\driftc\rho_0/M=m^\star$. The
\emph{running} fluctuation is controlled by the maximal inequality: applying
\Cref{lem:azuma} to the martingale $\bigl(\sum_{s\le t}w_{t,s}\inner{\rowhat}{\xi_s}\bigr)_{t\le T}$
(increments $\le2w_{t,s}M$, conditional variance $\le w_{t,s}^2\sigma_{\max}^2/d$),
\begin{align*}
   \max_{T_0\le t\le T}\Bigl|\inner{\rowhat}{\tilde x_t}-\E[\inner{\rowhat}{\tilde x_t}\mid\Fcal_0]\Bigr|
   \;\le\; 2M\sqrt{\tfrac{2S_T\log(2/\delta_{\mathrm{fail}})}{d}}
   \;\le\; \varepsilon M/2,
\end{align*}
on the good event, once $T$ is large enough that the noise scale
$\sigma_T\asymp e^S/\sqrt{Td}\le\varepsilon/4$ (equivalently
$T\ge T_0=O(\sigma_T^{-2}\varepsilon^2)$ raw steps). Dividing by
$\norm{\tilde x_t}_2\le M$ and combining, $\snorder_t\ge m^\star-\varepsilon$
for all $T_0\le t\le T$, i.e.\ $x_t\in C$.

\smallskip\noindent\textbf{(ii) Terminal concentration.} At $t=T$ the
two-sided form of \Cref{lem:signal_floor} pins
$\inner{\rowhat}{\tilde x_T}\in[\snet\driftc\rho_0-\nu,\ \snet M+\nu]$ with
$\nu=2M\sqrt{2S_T\log(2/\delta_{\mathrm{fail}})/d}=O(1/\sqrt{Td})$, and
$\norm{\tilde x_T}_2\in[\snet\driftc\rho_0-\nu,M]$ (the lower bound since
$\norm{\tilde x_T}_2\ge\inner{\rowhat}{\tilde x_T}$); hence $\snorder_T$ lies
within $O(\nu)=O(1/\sqrt{Td})$ of $m^\star$, which is
\Cref{eq:R3prime_concentration}. The verifier reads the on-sphere $x_T$,
but $\snorder_T$ is invariant under the positive rescaling
$x_T=\sqrt d\,\tilde x_T/\norm{\tilde x_T}$ (it is the same cosine), so the
retraction transfer (\Cref{lem:retraction_stability}) contributes only the
$O(1/d^{1.5})$ angular term, absorbed. This completes the proof.
\end{proof}

\begin{remark}[The ODE is a flagged stylization, NOT derivable from A1]
\label{rem:R3prime_ode}
It is tempting to summarise \Cref{thm:R3prime} by the scalar drift ODE
\begin{align}\label{eq:R3prime_ode}
   \dot m \;=\; \frac{\snet\,\driftc\,\rho_0}{M} \;-\; m
   \qquad\text{(informal; not used in any proof)},
\end{align}
whose stable fixed point $m=\snet\driftc\rho_0/M$ matches the target level
$m^\star$. We emphasise that \Cref{eq:R3prime_ode} is \emph{not derivable
from \Cref{ass:descent} alone}, and no result in this paper uses it. Two
additional stylizations beyond A1 would be required to make it rigorous,
neither of which we assume:
\begin{itemize}[leftmargin=2em,itemsep=0.2em]
\item \textbf{(AS) Asymptotic stationarity.} The loss-stratified descent
       process $(\descent_t,\loss_t)$ must be asymptotically stationary
       (e.g.\ ergodic, or $\alpha$-mixing with summable coefficients) for
       a time-averaged drift to converge to a state function of $m$ alone
       (cf.\ the order-parameter precedent of \cite{saad1995online} and
       the high-dimensional SGD scaling limits of
       \cite{benarous2022highdim}).
\item \textbf{(L$\lambda$) Lipschitz closure.} The conditional drift
       $\E[\snorder_t-\snorder_{t-1}\mid\Fcal_{t-1}]$ must be expressible
       as a Lipschitz function of $\snorder_{t-1}$ alone (a closability
       hypothesis), so that the finite-$d$ stochastic recursion has a
       well-posed ODE limit.
\end{itemize}
\Cref{thm:R3prime} deliberately avoids both: it is a finite-horizon
confinement statement, established by a maximal inequality over the
fixed horizon $T_{\max}$, and carries no hidden ODE or closability claim.
Because the running-average weights satisfy $\sum_t w_{T,t}=1$
(\Cref{lem:running_average}), the trajectory has no per-step
Foster--Lyapunov contraction rate (\Cref{rem:sgdm_reading}); attempting to
force a per-step rate is exactly what \Cref{eq:R3prime_ode} would smuggle
in, and we do not.
\end{remark}
